\documentclass{article}
\usepackage[utf8]{inputenc}
\usepackage{amsmath}
\usepackage{amssymb}
\usepackage{hyperref}
\usepackage{geometry}
\geometry{
 a4paper,
 margin=1in,
}
\begin{document}
\section{Introduction}
\textbf{Research question:} Do Google search indicators (aligned to W34 of month $M-1$ and W12 of month $M$) improve \emph{nowcasts} of German unemployment relative to autoregressive benchmarks, and does “early” W34 information deliver most of the gains?

\begin{itemize}
  \item Motivation: unemployment is released with a lag; nowcasting matters for policy and business decisions.
  \item Why Google Trends: high-frequency behavioral signal; prior evidence for labor markets.
  \item W34/W12 timing follows Askitas \& Zimmermann’s institutional logic for Germany (BA mid-month cutoffs and announcement cycle).
  \item Contributions: (i) bi-weekly alignment to the administrative window; (ii) early vs.\ late nowcast comparison; (iii) rigorous OOS evaluation on a short sample.
\end{itemize}



\section{Data}
\begin{itemize}
  \item \textbf{Target variable:} German unemployment rate (source, adjustment type, frequency, and seasonality note). Sample window: 2011-05 to 2023-01 (140 months). COVID-19 years (2020–21) serve as a stress-test period.
  \item \textbf{Google Trends:} 
  Weekly search intensities for five labor-market categories:
  \textbf{Inflow-type}:(“Arbeitsamt”, “Kündigung”, “Arbeitslosengeld”).
  \textbf{Outflow-type}: (“Jobbörse”, “Bewerbung”),
aggregated to bi-weekly W34prev and W12 components.  
  \item \textbf{Aggregation:} Weekly Google series averaged into bi-weekly components  (W34($M-1$) and W12($M$)) to match BA reporting cycles. Stitched overlapping Google windows (W1–W3) to consistent 2011–2023 series.
\end{itemize}





\section{Methodology}
\begin{itemize}
  \item \textbf{Pre-tests:} Augmented Dickey–Fuller (ADF) for stationarity; 
  \item Engle–Granger tests (trend = ‘c’ / ‘ct’).Evidence of cointegration between unemployment and arbeitsamtW34prev, jobboerseW34prev. Others: I(1) without cointegration.
  \item \textbf{Transformations:} Main specification in first differences ($\Delta$) to model short-run dynamics.
\end{itemize}




\subsubsection*{Path A — Cointegration / Error–Correction Approach}

If unemployment and one or more Google Trends indicators are found to be cointegrated,
a long–run equilibrium relationship can be established and modeled via an
Error–Correction Model (ECM). The ECM captures both short–run adjustments
and the speed of correction toward the long–run path:

\begin{equation}
\Delta y_t 
= \alpha 
+ \phi\,\Delta y_{t-1}
+ \beta_0\,\Delta x_t^{W34}
+ \beta_1\,\Delta x_{t-1}^{W34}
+ \gamma\, EC_{t-1}
+ D_m
+ \varepsilon_t,
\end{equation}

where $y_t$ denotes unemployment, $x_t^{W34}$ the ``early'' (weeks 34 of month $M-1$)
Google indicator, and $EC_{t-1}$ the lagged error–correction term obtained from the
cointegrating regression
\begin{equation}
y_t = a + b\,x_t + u_t.
\end{equation}

Seasonal dummies $D_m$ account for systematic month effects.  
Two ECM variants are estimated:

\begin{itemize}
  \item \textbf{Early ECM:} uses only the long–run relation and the $\Delta x_t^{W34}$ dynamics.
  \item \textbf{Late ECM:} augments the short–run equation with the latest
  available Google information from weeks 1–2 of month $M$ ($W12$):
  \[
  \Delta y_t = \alpha + \phi\,\Delta y_{t-1} 
  + \beta_0\,\Delta x_t^{W34} + \beta_1\,\Delta x_{t-1}^{W34}
  + \theta_0\,\Delta x_t^{W12} + \theta_1\,\Delta x_{t-1}^{W12}
  + \gamma\,EC_{t-1} + D_m + \varepsilon_t.
  \]
\end{itemize}

Out–of–sample one–step–ahead nowcasts are generated in an expanding–window design
(starting in May 2016, minimum training window = 36 months).  
Forecast accuracy is evaluated using RMSE and MAE.

\section{Nowcasting Design}
\begin{itemize}
  \item \textbf{Early nowcast:} Uses only 
W34(M−1) data — available sooner. (faster availability).
  \item \textbf{Late nowcast:} adds W12($M$) when available; measure incremental gain.
  \item Evaluation: expanding-window forecasts from 2016-05 onward (min 36 months training).
  \item Metrics: RMSE, MAE, directional accuracy; Diebold–Mariano tests.
\end{itemize}



\section{Results}
\begin{itemize}
  \item ADF \& Cointegration: Stationarity summary table; cointegrated pairs highlighted.
  \item Arbeitsamt ECM: strong long-run adjustment (γ ≈ –0.06).
  Jobbörse ECM: positive ΔW12 effect (short-run signal).
  \item \textbf{Out of sample}Report RMSE/MAE for early vs. late ECMs. Late ECM shows marginal but consistent gains, especially for outflow searches.
\end{itemize}




\section{Robustness and Extensions (still checking feasibility)}
\begin{itemize}
  \item Seasonality treatment: month dummies vs. Δ12 vs. YoY changes.
  \item AAlternative lag orders (1–12).
  \item Keyword grouping or PCA composites.
  \item COVID subperiod robustness.
\end{itemize}

\section{Conclusion}
\begin{itemize}
  \item Google search indicators contain predictive content for German unemployment.
  \item When cointegration holds, ECM captures long-run alignment; when not, ARX suffices.
  \item “Late” information (W12) improves short-run adjustment—especially for job-search (outflow) queries.
  \item Implication: Real-time monitoring via search data can complement official labor statistics.
\end{itemize}




\end{document}